\documentclass[11pt]{article}
\usepackage[T1]{fontenc}
\usepackage[margin=1in]{geometry}
\usepackage{amsmath,amssymb,amsthm,microtype}
\usepackage[colorlinks=true,allcolors=blue]{hyperref}
\newtheorem{theorem}{Theorem}[section]
\newtheorem{lemma}[theorem]{Lemma}
\newcommand{\F}{\mathbb F_2}
\newcommand{\Tr}{\operatorname{Tr}}
\title{Satisfying parity constraints on a Boolean truth table with a small circuit}
\author{Research manuscript}\date{22 November 2026}
\begin{document}\maketitle
\begin{abstract}
Given $r$ consistent parity constraints on the entries of an $N=2^n$-bit truth table, how many gates suffice to compute a solution? At $r=\Theta(N)$ the answer in the worst case is $(r+o(N))/n$. For some absolute $K$ and all $n\log n(\log\log n)^K\le r\le N$, the answer is $(1+o(1))r/\log r$. The second construction provides one common distribution of small circuits whose outputs look uniform even under $r$ adaptive parity queries. A shared lookup across many blocks determines the leading constant for bounded-error parity queries to Boolean MCSP in this range. At a fixed fraction $c$ of Shannon's size threshold a companion proof gives the leading constant even for two-sided error. The largest affine spaces consisting wholly of accepted or wholly of rejected tables have dimensions $(c+o(1))N$ and $(1-c+o(1))N$, respectively.
\end{abstract}

\section{What a parity constraint specifies}
Each row of an $r$ by $N$ binary matrix $A$ describes a parity query on a truth table $x$. We assume the rows are linearly independent. A right-hand side $y\in\F^r$ specifies the desired answers. A Boolean circuit for $x$ takes the $n$-bit address of an entry and returns its value. We count binary AND and OR and unary NOT gates with free constants and unrestricted fanout; arbitrary binary gates obey the same asymptotics. Write
\[
 L(n,r)=\max_{\operatorname{rank}A=r,\ y\in\F^r}
             \min\{ C(x): Ax=y\}.
\]
Here $C(x)$ is the minimum number of gates in such an address circuit. The linear equations can involve any subsets of the truth table. All logarithms are base two.
\begin{theorem}\label{th:sharp}
If $r=\Theta(N)$, then $L(n,r)=(r+o(N))/n$. For a sufficiently large absolute $C$ and all $n^C\le r\le N$, one has $L(n,r)=\Theta(r/\log r)$.
\end{theorem}
Specifying individual entries is a special case of the problem. The classical theory of partial Boolean functions, including Chashkin's synthesis theorem \cite{chash}, controls that special case. A circuit satisfying specified entries can take arbitrary values elsewhere. For general parity constraints those values may change the answers. Hoza and Lv \cite[Corollary 32]{hl} previously used successive rank increments in blocks in their hitting set for general consistent parity constraints. Here each solution must itself have a small Boolean address circuit.

A rooted binary circuit with $s$ gates has only $s+1$ predecessor slots beyond those that introduce a gate. Its description uses at most $(s+1)\log(n+s+2)+O(s+\log n)$ bits, including gate types and traversal structure. Since a fixed rank-$r$ matrix has $2^r$ possible answers, this count proves the lower bounds in the theorem. We next give two constructions for the upper bounds.

\section{The leading constant for a dense set of constraints}
Split the $N$ positions into consecutive blocks of $b$ positions, where $b$ is a power of two between $n^2$ and $2n^2$. For each block $J$, take the span of its columns modulo the span of all earlier blocks. Write $d_J$ for this dimension. Then $0\le d_J\le b$ and $\sum_J d_J=r$.

For each $0\le d\le b$, fix a list of $O(b^2+1)$ binary $b$ by $d$ matrices with the following property. For every rank-$d$ matrix $L:\F^b\to\F^d$, some list element $G$ has $LG$ invertible. For $d>0$ a uniformly random $G$ works for fixed $L$ with probability $\prod_{i=1}^d(1-2^{-i})$, bounded below by an absolute positive constant. There are at most $2^{bd}$ possibilities for $L$, proving the list claim by the union bound; $d=0$ is immediate.

On block $J$ choose a list element $G_J$ for the quotient map in any basis. Restrict $x|_J$ to the form $G_J u_J$ for a message $u_J$ of $d_J$ bits. By induction on the blocks the images of all restricted block generators through $A$ span precisely the same space as all original columns seen so far. Hence the $r$ message bits together can realize every syndrome $y$. The choice of generators depends on $A$ and can be shared across all its syndromes.

It remains to charge for all data a circuit needs at an address. A directory indexed by the block stores its rank, list index and ordinal among blocks of the same rank. These require $O(n)$ output bits on $\log(N/b)$ address bits. Implementing each by Lupanov's ordinary synthesis bound \cite{lup} costs $O((N/b)/\log(N/b))$ gates, for total $O(N/b)=o(N/n)$. The common generator lists, their selectors and their matrix-vector products need only $\operatorname{poly}(b)$ gates, which is also $o(N/n)$.

For each rank $d$ with $t_d$ blocks, store $d$ ordinary Boolean tables of length $t_d$, one per column of their messages in block-ordinal order. All these tables together contain exactly $r$ bits. Use at most $O(b^2)$ such tables, evaluated in parallel, and select the returned $d$ bits according to the directory. For clarity about the leading constant, split any table of length $t_d$ into dyadic blocks in its binary prefix decomposition. Across the polynomially many tables, pieces of length below $2^{n-\sqrt n}$ cost $N\,2^{-\sqrt n}\operatorname{poly}(n)=o(N/n)$ even by direct synthesis. A larger piece of length $2^m$ has $m\ge n-\sqrt n$ and costs at most $(1+o(1))2^m/n$, uniformly by Lupanov. Prefix comparisons and multiplexing over the polynomially many pieces have polynomial cost. Thus the total message circuit costs $(r+o(N))/n$. This proves the dense upper bound, including the overhead of locating each message.

\section{A field construction for fewer constraints}
The next proof has two uses. It reaches smaller $r$, and its distribution is independent of the matrix being queried. A related object in the literature is a \emph{lossless rank condenser}, a family of linear maps among which one preserves the dimension of any prescribed low-dimensional subspace. We include the rank argument and the Boolean address-circuit accounting. See Forbes and Guruswami \cite{fg} for algebraic rank condensers and their precedents. Cheraghchi \cite[Corollary 2.23]{cher} gives a binary linear lossless condenser based on Guruswami, Umans and Vadhan. Transposing a seeded binary linear condenser and selecting a uniform message gives the same general mechanism for affine tests: every good seed makes a prescribed rank-$r$ system jointly uniform. Our estimate includes the gate budget for a fixed output word.

Regard a vector $a\in\F^N$ as the coefficient list of $P_a(X)=\sum_{i=0}^{N-1}a_iX^i$. Set
\[
 \ell=n^2,\qquad h=\ell-4\lceil n+\log(r+1)+2\rceil,
 \qquad d=\lceil r/h\rceil,\qquad F=\mathbb F_{2^\ell}.
\]
Take $n$ sufficiently large so $h>0$. Pick $v$ uniformly in $F$, treating $v=0$ as a bad seed and defining $v^0=1$. Consider the binary-linear map taking $a$ to $(P_a(v^{2j-1}))_{1\le j\le d}\in F^d$.
\begin{lemma}\label{lem:rank}
For every fixed rank-$r$ binary row space $U$, the displayed map is injective on $U$ except with probability $O((1+rNd\,2^h)/2^\ell)=2^{-\Omega(n)}$ over $v$.
\end{lemma}
\begin{proof}
Choose an echelon basis of $U$, giving polynomials $P_1,\ldots,P_r$ with distinct leading degrees. From the positive integers $(2j-1)2^a$ for $1\le j\le d$ and $0\le a<h$ choose any $r$, in increasing order. They are distinct by their odd parts. In the polynomial determinant with entries $P_i(Z^{e_j})$, pair the increasing leading degrees with the increasing exponents $e_j$. This permutation gives the unique term of largest degree by the strict rearrangement inequality. Every lower monomial in the product has strictly smaller degree. The determinant is nonzero even in characteristic two. Its degree is at most $rN(2d)2^h$. Counting $v=0$ as an additional failure and counting the other roots yields the failure bound by counting roots. A binary dependence among the vectors $(P_i(v^{2j-1}))_j$ would remain a dependence after raising every component to any of the $2^a$ powers, since the polynomials have binary coefficients. It would therefore annihilate the invertible specialized determinant.
\end{proof}

Independently choose uniform $w_1,\ldots,w_d$ in $F$ and form the Boolean table
\[
 x_i=\Tr_{F/\F}\!\left(\sum_{j=1}^d w_j(v^i)^{\,2j-1}\right).
\]
For a fixed query row $a$, its answer is $\Tr(\sum_j w_j P_a(v^{2j-1}))$. Thus if $v$ is good in Lemma~\ref{lem:rank}, the $r$ answers for the rows of $A$ are uniform jointly as $w$ varies. We use here the nondegeneracy of the trace pairing. In particular any consistent assignment is attained. Unconditionally, the statistical distance of $Ax$ from uniform is $o(1)$ for every fixed rank-$r$ matrix.

\paragraph{Cost at a binary address.}
Compute $v^i$ by successively multiplying by a fixed $v^{2^a}$ when address bit $a$ is one. This costs $\operatorname{poly}(n)$ since $\ell=n^2$, and all its field operations can be done by schoolbook binary operations. A fixed degree-$O(d)$ polynomial over $F$ has an evaluation circuit of size
\[
                    O(d\ell/\log(2d)+\sqrt d\,\operatorname{poly}(\ell)).
\]
Here are the details of the factor $\ell$ in its leading term. Write the coefficients in a square array of side $m=\lceil\sqrt{2d+1}\rceil$. Precompute $m$ small powers of the variable and $m$ powers of its $m$th power. Binary extension field multiplication has bilinear rank $R=O(\ell)$ \cite{rank}. For each of the $R$ terms, take the corresponding input linear form on all $m$ small powers. One fixed $m$ by $m$ binary map forms that term's $m$ inner sums of hardwired coefficients times these bits. Each binary map costs $O(m^2/\log(2m))$ XORs by grouping inputs into blocks of $\lfloor\log(m)/2\rfloor$ bits, precomputing subset sums and combining one result per block per output. Recombining all terms gives the block polynomials, which are multiplied by the powers of the $m$th power and summed. The initial and final linear transforms and the $O(m)$ ordinary field multiplications cost $O(m\operatorname{poly}(\ell))$. Finally the trace is a binary linear form. Consequently when $r\ge n^C$ for sufficiently large fixed $C$, the entire cost is $O(r/\log r)$, uniformly. This proves Theorem~\ref{th:sharp}.

\section{Consequences for Boolean MCSP}
Let $A_s=\{x:C(x)\le s\}$. At $s\sim cN/n$ for $0<c<1$, the greatest dimension of an affine subspace contained in $A_s$ is $(c+o(1))N$. Counting gives the upper bound. Conversely all tables supported on a prefix of $(c-\delta)N$ entries can be synthesized with $(c-\delta+o(1))N/n$ gates by the same dyadic lookup calculation, for each fixed $\delta>0$.

The greatest dimension of an affine subspace in its complement is $(1-c+o(1))N$. The dense syndrome theorem forces every affine space defined by at most $(c-\delta)N$ independent queries to hit $A_s$. Conversely a fixed full-rank matrix with $(c+\delta)N$ rows has more syndromes than $|A_s|$; some coset misses $A_s$. Let $\delta\downarrow0$ in the two bounds. Thus the least rank of an affine certificate that a table is hard, minimized over such certificates and tables, has order $(c+o(1))N$.

More generally consider a randomized decision tree whose queries are parities of arbitrary subsets of the truth-table entries, with unbounded computation. At $n^{C'}\le s\le(1-\epsilon)N/n$, for sufficiently large absolute $C'$ and any fixed $\epsilon>0$, its bounded-error query cost is $\Theta_\epsilon(s\log s)$, even when the choice of each parity depends on earlier answers.

For the lower bound take $r=\lfloor c_0s\log s\rfloor$ for small $c_0>0$. Our common field construction has every output in $A_s$. Fix a deterministic depth-$r$ tree. Each reachable leaf prescribes a fixed consistent system of at most $r$ linear equations. With probability at least $1-\delta$ over $v$, where $\delta=2^{-\Omega(n)}$, the projection is surjective on the independent rows of this system. To see this one can extend their span to rank $r$ before applying Lemma~\ref{lem:rank}. At every such seed the probability, over the message, of that leaf equals its exact probability under a uniform input. Thus each leaf receives at least $(1-\delta)$ times its uniform probability under the common easy law. Summing the smaller probability of the two laws at each leaf gives an overlap of at least $1-\delta$ over all leaves, regardless of their number. Averaging over tree coins gives the same conclusion for a randomized tester. A uniform table lies outside $A_s$ with probability $1-o(1)$ by counting, so a correct tester must make more queries. For the upper bound draw $\lceil\log |A_s|\rceil+2$ independent random parity rows and accept if any easy table has the observed syndrome. Every member passes, and each fixed nonmember collides with any given member with probability $2^{-q}$ for $q$ rows.

The same argument gives a generator for every depth-$r$ parity decision tree. With $v$ and $w$ drawn as above, its seed has $\ell+d\ell=(1+o(1))r$ bits for $r\ge n^C$ with $C$ large enough. The statistical error of the full transcript is exponentially small in $n=\log N$. A seed of $r-O(1)$ bits is necessary at constant error because an observer can read $r$ individual bits and check membership in their projected support. Standard small-bias generators already fool parity decision trees with seed length of optimal asymptotic order, as explained by Hoza and Klivans \cite[Section~1.5.4]{hk}. Here every output is also an address function with $O(r/\log r)$ Boolean gates.

For adaptive randomized parity queries with perfect completeness, the same order holds, and at $s\sim cN/n$ the cost is $(c+o(1))N$. A deterministic rejecting transcript describes a nonempty affine space of codimension at most its query length. Perfect completeness disallows any such transcript intersecting $A_s$. This reasoning applies after fixing the random coins outside a common null set because $A_s$ is finite. The field construction gives the order lower bound and the dense construction gives the leading constant. The same random linear sketch supplies the upper bounds.
A companion manuscript, \emph{Small Boolean truth tables under adaptive parity tests}, sharpens the coefficient at dense thresholds. A random affine permutation spreads a basis of the queried columns across the blocks. Independent small matrix choices retain all but a vanishing fraction of their rank and a second code restores the missing directions. The resulting common distribution proves the leading $(c+o(1))N$ lower bound even for two-sided error, with an application to unrestricted outer circuits.

\section{Prescribing linear information before testing circuit size}
Suppose $s\sim cN/n$, with fixed $0<c<1$. A rank-$r$ matrix $A$ specifies $r$ parity checks on a table. Its syndrome is the vector $Ax$ of answers. For a fixed syndrome $y$ we may treat $Ax=y$ as known in advance and ask what remains of MCSP on that affine space. Neither the matrix nor the choice of coordinates on its solution space need be easy to compute.

\begin{theorem}
Fix $0<\rho<c<1$ and $r=\lfloor\rho N\rfloor$. Uniformly for every rank-$r$ matrix $A$, each fiber $Ax=y$ contains an affine space of easy tables of dimension at least $(c-\rho-o(1))N$. One may use the same direction space for all $y$ for this $A$. If the syndrome is uniform, then with probability $1-o(1)$ the following statements hold:
\begin{enumerate}
\item The largest affine space of easy tables in the fiber has dimension $(c-\rho+o(1))N$, and the largest affine space of hard tables there has dimension $(1-c+o(1))N$.
\item The number of easy tables in the fiber has logarithm $(c-\rho+o(1))N$. The bounded-error parity query complexity of MCSP restricted to this fiber, and the logarithm of its bounded-error minimum parity tree size measured in leaves, are both $(c-\rho+o(1))N$.
\end{enumerate}
If instead $c<\rho\le1$, the number of attainable syndromes of easy tables has logarithm $(c+o(1))N$, for every $A$. A uniform syndrome is attainable with probability $2^{-(\rho-c+o(1))N}$.
\end{theorem}
\begin{proof}
For any fixed $0<\epsilon<c-\rho$, extend the $r$ rows to a linearly independent list of $r'=\lfloor(c-\epsilon)N\rfloor$ rows, denoted by $A'$. The block construction of Section~2 yields one linear image $U$ of dimension $r'$ that projects surjectively under $A'$, with every word computable in $(r'+o(N))/n\le s$ gates. The common direction $U\cap\ker A$ has dimension $r'-r$. Intersection of $U$ with each fiber is a coset of this space. Sending $\epsilon$ slowly to zero gives the first statement.

Let $M_y=|A_s\cap\{Ax=y\}|$. By the circuit census,
$\mathbb E_y M_y=|A_s|2^{-r}\le 2^{(c-\rho+o(1))N}$.
Markov's inequality and the preceding lower bound then give $\log M_y=(c-\rho+o(1))N$ with high probability. The same upper bound applies to the dimension of any easy flat in that fiber. Within such a fiber, any rank-$k$ internal parity check splits it into $2^k$ cosets. If $2^k>M_y$, one of these cosets consists entirely of hard tables. Take $k=\lceil\log(M_y+1)\rceil$. Its dimension is $N-r-k=(1-c-o(1))N$. The global hard-flat upper bound in the preceding section applies to all such cosets.

For the query lower bound use the common random linear encoder of the companion manuscript \emph{Small Boolean truth tables under adaptive parity tests}, Theorem~2, at rank $r'$. The ensemble is independent of all queried rows, its failure on any fixed rank-$r'$ system is $o(1)$, and every message maps to an easy table. Given $A,y$, draw an encoder conditioned on full projection onto $A$, then draw a uniform message among those mapping into the fiber. For any fixed additional $q\le r'-r$ tests independent modulo $A$, extend the joint list to rank $r'$. With conditional probability $1-o(1)$ the encoder has full projection on the list, since the conditioning event has probability $1-o(1)$. For each such encoder all $2^q$ additional syndromes are equiprobable at the prescribed $y$. Each affine cell of codimension $q$ inside the fiber thus has at least $(1-o(1))2^{-q}$ mass under this law. The error bound is uniform over cells. A parity tree with at most $2^{r'-r-\epsilon' N}$ leaves assigns at most $2^{-\epsilon'N}$ of the fiber's uniform mass to leaves of codimension above $r'-r$. Summing the smaller masses over the other leaves shows that the two transcript laws have total variation $o(1)$, even after averaging over independent tree coins. In a typical fiber the uniform hard probability tends to one by the bound on $M_y$. Correct bounded-error recognition therefore requires larger trees. A depth-$q$ tree has at most $2^q$ leaves, and $\lceil\log M_y\rceil+2$ random parity checks suffice by the fingerprint argument. Taking the fixed margins to zero proves the two query conclusions in probability.

Finally, when $\rho>c$, the circuit census bounds the number of syndromes by $2^{(c+o(1))N}$. Restrict $A$ to any $\lfloor(c-\epsilon)N\rfloor$ of its rows. Every partial syndrome is attained by an easy table by Section~2, giving at least $2^{(c-\epsilon)N-O(1)}$ distinct complete syndromes. This proves the last assertion.
\end{proof}

The typical-syndrome qualification matters. For instance an affine space of dimension less than $(c-\epsilon)N$ supported on an initial segment of addresses can consist entirely of easy tables, by the lookup bound of Section~2. If $c+\rho<1$, however, the fraction of uniform tables that are easy tends to zero in \emph{every} fiber by the global count. In this regime the parity query lower bound applies in every fiber.

\section{The leading constant for a polynomial or larger number of constraints}
The field calculation becomes more efficient if its decoder is shared among many blocks. For a truth table on $n$ bits let $N=2^n$. Throughout this section $n^{30}\le r\le N$. The exponent $30$ keeps the secondary costs explicit.

\begin{theorem}\label{th:sparse-sharp}
Uniformly in this range, $L(n,r)=(1+o(1))r/\log r$. There is one distribution $\mu_r$ supported on tables of size at most $(1+o(1))r/\log r$ such that, for every rank-$p$ matrix $A$ with $p\le r$ and every $y\in\F^p$,
\[
                         \Pr_{\mu_r}[Ax=y]\ge(1-o(1))2^{-p}.
\]
\end{theorem}
\begin{proof}
It suffices to consider $p=r$ because a row basis can be extended. The lower bound for $L(n,r)$ follows from rooted circuit counting: every circuit of size below $(1-\epsilon)r/\log r$ has a description of fewer than $r$ bits, with a fixed linear margin, for any fixed $\epsilon>0$ and sufficiently large $n$. Consider a fixed rank-$r$ matrix and mark a basis of $r$ of its columns.

\paragraph{Distributing the marked columns.}
Let $B$ be the largest power of two at most $r/n^{20}$, and set $\lambda=r/B$, so $n^{20}\le\lambda<2n^{20}$. Split addresses into $B$ equal blocks after a uniform invertible affine permutation on $\F^n$. The number $p_J$ of marks in each block has mean $\lambda$ and variance at most $\lambda$, by transitivity on ordered pairs. Call a block crowded if $p_J>\lambda+\lambda^{3/4}$. Its expected number of marks charged on this event is at most
\[
 \lambda\,\frac{\lambda}{\lambda^{3/2}}
       +\frac{\lambda}{\lambda^{3/4}}=O(\sqrt\lambda).
\]
Indeed $\mathbb E[(p_J-\lambda)1_{\rm crowded}]\le
\mathbb E[(p_J-\lambda)^2]/\lambda^{3/4}$. Summing over blocks bounds the expected number of crowded marks by $O(r/n^{10})$.

\paragraph{Coding in each block.}
Use Lemma~\ref{lem:rank} on the $H=N/B$ positions of each block, with local target $q=\lceil\lambda+\lambda^{3/4}\rceil$, working field degree $\ell=n^2$, $h=\ell-4\lceil n+\log(q+1)+2\rceil$ and $d=\lceil q/h\rceil$. This gives a binary linear local encoder with selector $v$ of $\ell$ bits, message $(w_j)_{j\le d}$ of $d\ell$ bits, and success probability at least $1-N^{-2}$ for every fixed test of rank at most $\min(q,H)$. If $q>H$, apply the determinant proof directly to the actual number of independent rows, using that there are at least that many distinct evaluation exponents. The number of selector and message bits in each block is
\[
                   W=(1+d)\ell=(1+o(1))\lambda .
\]
The logarithm in the error bound uses $n$ as an upper bound for $\log H$, and $q=O(n^{20})$. Selectors and messages are independent between blocks.

Process the blocks while retaining an independent set. Initially retain all marked columns of $A$. Discard all marks in a crowded block. For any other block, the $p_J$ original marks remain independent of all the other retained columns, including replacements from preceding blocks. Take $p_J$ forms that annihilate those other columns and give the coordinate basis on the current marks. On the $H$ columns here they define a fixed rank-$p_J$ test, independent of the next selector. If the local code preserves its rank, replace the $p_J$ marks by $p_J$ coded columns independent modulo the other retained columns. Otherwise discard its marks. The conditional expected extra loss per such block is at most $p_J/N^2$. Thus the overall expected rank deficit is $O(r/n^{10})$. By Markov's inequality the deficit is at most $t=\lceil r/n^6\rceil$ with probability $1-O(n^{-4})$.

\paragraph{Recovering the remaining rank.}
Add a second, independent table, at original address $i$, with bit
$\operatorname{Tr}\sum_{j\le D} w'_j (v'_j)^i$, where $D=\lceil 2t/n\rceil+1$ and all the primed elements are independent uniform in $\mathbb F_{2^{2n}}$. Condition on the local selectors and permutation when the deficit is at most $t$. Each of the at most $2^t-1$ nonzero missing combinations of the rows of $A$ is a nonzero coefficient polynomial of degree below $N$. The probability that it vanishes at all of the $D$ independent points is at most $2^{-nD}$. Except with probability $2^{t-nD}$, the combined encoder therefore has rank $r$, by nondegeneracy of the trace pairing. Independent uniform messages now make every right-hand side equally likely.

\paragraph{Reading all the bits for one block.}
Given an address, the circuit computes its affine image and then retrieves the $W$ selector and message bits stored on its block. We give the lookup cost explicitly because the number of output wires grows with $n$. Put $P=BW=(1+o(1))r$. There is a power of two $M$ between $(\log P)^2$ and $2(\log P)^2$, with $M\mid B$ for large $n$. The low address has $M$ possible values. Partition them into groups of $g=\lfloor\log P-6\log\log P\rfloor$ values, with a possible shorter final group. Construct all $M$ low-address indicators and all $B/M$ high-address indicators using $O(M+B/M+\log B)$ gates by recursive extension of prefixes. For each low group precompute the OR of every subset of its indicators in $O(2^g)$ gates, by extending subsets one bit at a time. These shared costs are $o(P/\log P)$.

For each of the $B/M$ high addresses and each output bit, choose the desired subset OR in each low group, combine the selected terms, and gate by the high-address indicator. Then take the union over high addresses. The cost is at most
\[
           \frac{WB}{M}\bigl(\lceil M/g\rceil+2\bigr)
                      =(1+o(1))\,\frac P{\log P}.
\]
This is the grouping argument in Lupanov's synthesis method, with the groups reused across all outputs. It gives the leading constant also when $\log B$ and $\log P$ differ by a fixed factor.

\paragraph{One decoder for every block.}
The same hardware evaluates the selected field polynomial on the position within the block, using the retrieved $v,w_j$ as inputs. Repeated squaring and conditional multiplication compute the field power in $O(n\ell^2)$ gates. Horner evaluation costs $O(d\ell^2)$ further gates using schoolbook extension field multiplication. Altogether these gates cost $O(n^{22}+n^5)=o(r/\log r)$; the affine permutation also fits. The second table uses $O(Dn^3)=O(r/n^4+n^3)=o(r/\log r)$ gates by the same schoolbook estimate. One final XOR completes the evaluation. The endpoint estimates are uniform for $n^{30}\le r\le N$.
\end{proof}

Consequently in the range $n^{30}\le s\le(1-\epsilon)N/n$, for each fixed $\epsilon>0$, bounded-error adaptive parity query complexity and the logarithm of minimum bounded-error parity-tree size are both $(1+o(1))s\log s$. In the lower bound use $r=\lfloor (1-\gamma)s\log s\rfloor$ for each fixed $\gamma>0$. The easy law gives relative lower bounds simultaneously at the level of all leaf probabilities. A tree with at most $2^{(1-\gamma)r}$ leaves has vanishing uniform probability on leaves of rank greater than $r$. At every other leaf the two distributions have total overlap at least $1-o(1)$ after summing the smaller masses. A uniform input is hard with probability $1-o(1)$. The random parity sketch in the preceding section attains the matching upper bound by rooted circuit counting. The same conclusion holds with a larger promised no threshold up to $(1-\epsilon)N/n$.

\section{Bringing the gate budget near the address length}
We sharpen Theorem~\ref{th:sparse-sharp} at its lower endpoint. This section changes the representation of a row polynomial. It makes the cost of evaluating a field element depend almost linearly on the address length.

\begin{theorem}\label{th:near}
There is an absolute $K$ such that Theorem~\ref{th:sparse-sharp} holds whenever
\[
             n\log n\,(\log\log n)^K\ \le r\le N .
\]
In particular the leading constant for bounded-error parity queries and log leaf count is one in units of $s\log s$ for $2n(\log\log n)^K\le s\le(1-\epsilon)N/n$, for every fixed $\epsilon>0$.
\end{theorem}
The extra factor $2$ at the lower endpoint is only used to choose a fixed margin for the lower bound. The simple odd-weight equation already forces an $n$-variable solution to depend on all its variables. Every function ignoring one variable has an even number of ones. Thus the raw address length is a genuine scale in the interpolation problem.

\subsection{Rows as polynomials in address digits}
We need only consider $r\le n^{30}$ by the previous theorem. Set $L=\lceil\log n\rceil$, $J=\lceil\log(L+1)\rceil$, $\ell=LJ$, and $h=\ell-120L$, taking $n$ sufficiently large. Write each address as $m=\lceil n/L\rceil$ digits $i_g$ of $L$ bits, padding the final digit with zeros. For a binary row $a$ on the addresses define
\[
              P_a(Z_1,\ldots,Z_m)=\sum_i a_i\prod_{g=1}^m Z_g^{i_g}.
\]
The row is recoverable from the coefficients. Any subset of the addresses can be used by setting the other coefficients to zero. For target rank $q\le n^{30}$ let $d=\lceil q/h\rceil$ and take independent nonzero $v_g$ in $F=\mathbb F_{2^\ell}$. Evaluate at the $d$ points with coordinates $v_g^{2j-1}$ for $1\le j\le d$.

\begin{lemma}\label{lem:digits}
For every fixed binary row space of rank at most $q$ these evaluations are jointly injective, except with probability at most $n^{-40}$. The corresponding transposed encoder with uniform messages $w_j\in F$ has value
\[
      x_i=\operatorname{Tr}\sum_{j=1}^d w_j Y_i^{\,2j-1},
                \qquad Y_i=\prod_{g=1}^m v_g^{i_g}.
\]
Its selector uses $m\ell=O(nJ)$ bits and its message uses $d\ell\le(1+O(1/J))q+O(\ell)$ bits. Given its selector, message and physical address as inputs, the output is computed by a single Boolean decoder of size $O((n+q)J^{K_0}+\ell^3)$ for some absolute $K_0$.
\end{lemma}
\begin{proof}
Choose an echelon basis with distinct leading monomials in lexicographic order. Use the determinant argument from Lemma~\ref{lem:rank}, raising each coordinate of the indeterminate tuple to its indicated exponent $e=(2j-1)2^a$ for $a<h$. Among the determinant products, pairing increasing leading exponent vectors with increasing $e$ gives the unique largest lexicographic vector. Swapping an inversion improves the vector by a positive integer times the lexicographically positive difference of the two distinct vectors. Replacing any leading monomial by a lower monomial reduces the sum in lexicographic order. Thus the multivariate determinant is nonzero. For a rank below $q$ use as many distinct exponents as the actual row rank.

Its total degree is at most $q\,m\,2^L(2d)2^h$. The Schwartz--Zippel bound on $(F^\times)^m$ puts the failure probability below $n^{-40}$, since $m\le n$, $2^L\le2n$, $d\le q+1$ and $\ell-h=120L$. Frobenius powers preserve a putative binary dependence. Hence a nonzero determinant forbids one, and the trace-pairing calculation gives the displayed encoder.

Fix a primitive field element $g_0$ and store each nonzero selector $v_g$ by its discrete logarithm $z_g$. The selector is sampled in that representation. At an address compute $\sum_g i_g z_g$ modulo $2^\ell-1$ and exponentiate $g_0$ by this one $\ell$-bit exponent to get $Y_i$. Products of $\ell$-bit integers and products of binary polynomials of length $\ell$ take $O(\ell(\log\ell)^{O(1)})$ Boolean gates \cite{intmult,polyprod}. Reduction of an integer product modulo $2^\ell-1$ needs only a constant number of $\ell$-bit additions and a conditional subtraction. Reducing a binary polynomial product modulo the fixed field polynomial also has quasi-linear size: reverse its leading coefficients, multiply by the precomputed truncated reciprocal of the reversed modulus to get the quotient, and subtract its product with the modulus. Thus computing the modular sum costs $m\ell(\log\ell)^{O(1)}=nJ^{O(1)}$. Ordinary repeated squaring for the one exponent costs $O(\ell^3)$. Horner evaluation in $Y_i^2$ followed by multiplication by $Y_i$ and trace costs $O(d\ell(\log\ell)^{O(1)})$. This proves the stated bound, including the cost when the coefficients and the selector are retrieved from a block.
\end{proof}

\subsection{Sharing the digit decoder}
Take $B$ the largest power of two at most $r/(nJ^4)$ and $\lambda=r/B$. Assign each physical address to a block using a pairwise independent affine binary hash with $\log B$ output bits. Only the block number is hashed. The local decoder uses the original physical address, so it needs neither a permutation nor an inverse. For any fixed pivot basis, the mark count in a block has mean $\lambda$ and variance at most $\lambda$. With $q=\lceil\lambda+\lambda^{3/4}\rceil$, use Lemma~\ref{lem:digits} with independent selectors and messages on each block. Its selector contributes $O(nJ)=o(\lambda)$ bits, and local messages have $(1+o(1))\lambda$ bits. The charging and pivot-replacement proof of Theorem~\ref{th:sparse-sharp} yields expected rank deficit $O(r/\sqrt\lambda)$, since the additional expected loss is at most $rn^{-40}$. With probability $1-O(\lambda^{-1/4})$ the deficit is at most $t=\lceil r/\lambda^{1/4}\rceil$. One further independent encoder of Lemma~\ref{lem:digits} with target $t$ on the entire physical table restores full rank except with probability $n^{-40}$.

The lookup table for the local messages and selectors has $P=(1+o(1))r$ bits. In the grouping proof take $M$ the least power of two above $\log P\log\log P$ instead of $(\log P)^2$. This still gives the stated leading constant because $M/\log P\to\infty$. For $r\ge nLJ^K$ and $K$ sufficiently large, $B\ge M$. Indeed $\log P=O(L)$ in the present range and $B\ge LJ^{K-4}/2$. The two shared decoders cost
\[
  O\bigl((n+\lambda+t)J^{K_0}+\ell^3\bigr)=o(r/\log r)
\]
when $K>K_0+6$: here $\lambda\le2nJ^4$, $\log r=O(L)$ and $t\le 1+r/\lambda^{1/4}$. The hash costs $O(n\log B)$ gates. For $r\le nL^3$ this is $O(nJ)$, and for larger $r$ it is $O(nL)=o(r/\log r)$. Thus the cost, full-rank probability and counting lower bound still hold. The same leaf argument proves the parity-query statement.

\subsection{A dimension reduction specific to recognition}
For \emph{exact recognition}, the parity lower bound extends to every $s\to\infty$ with $s\le(1-\epsilon)N/n$. It does so even when the interpolation construction on the original $n$-variable domain is too costly. Let
\[
                  m=\min\{n,\lceil 2\log(s+2)\rceil\}.
\]
Embed an $m$-variable table by making it independent of the other $n-m$ inputs. Its minimum circuit size stays exactly the same: a circuit on $m$ inputs extends directly, and fixing the unused variables in an $n$-input circuit gives the reverse inequality. One parity of the extended table is one parity of the original table, by collecting repetitions modulo two. Since $m=O(\log s)$, eventually $s\ge m^{30}$. When $m<n$, $2^m/m\gg s$; when $m=n$ the assumed margin suffices. Applying the preceding lower bound in dimension $m$ proves
\[
 q_{\rm parity}\ \ge\ (1-o(1))s\log s,\qquad
 \log_2(\text{parity-tree leaf count})\ \ge\ (1-o(1))s\log s.
\]
The fingerprint upper bound on the original domain is $(s+1)\log(n+s+2)+O(s+\log n)$. Consequently the leading constant is sharp throughout $\log(n+s)\sim\log s$, including all fixed positive linear gate budgets in $n$. For the genuine interpolation problem one must still meet the equations on their original domain, so the raw-address argument in Theorem~\ref{th:near} has independent content. A gap between the yes and no size thresholds also requires that the smaller domain still supports the promised no inputs.

\subsection{A bound independent of the number of address variables}
Even if there are arbitrarily many more raw variables than gates, exact MCSP has bounded-error parity query complexity $\Theta(s\log s)$ for $s\to\infty$ and $s\le(1-\epsilon)N/n$. The lower bound is the previous restriction. The following upper bound also has perfect completeness and uses one nonadaptive batch. The argument relies on the fact that a rooted $s$-gate circuit uses at most $k=s+1$ essential input variables.

Hash all raw variable names independently into $M=100(k+1)^2$ buckets. A fixed list of at most $k+1$ names hashes injectively except with probability below $1/100$. We first check whether the number of buckets containing an essential variable exceeds $k$. For each bucket $b$ form a binary vector consisting of all derivatives $f(x)+f(x\mathbin{\oplus} e_i)$, for $i$ in the bucket and all addresses $x$. It is zero precisely for an inactive bucket. Independently apply a random binary linear map to turn each of these vectors into a field symbol $V_b$ in $\mathbb F_{2^j}$ for $j=\lceil\log_2(100(k+1))\rceil$. On active buckets the symbols are uniform and independent. All binary coordinates of $V$ are linear functionals of the original table.

Choose $q=\lceil\log_2|\mathcal S|\rceil+7$ random binary parity rows on the concatenated vector $V$, where $\mathcal S$ consists of vectors with at most $k$ nonzero field symbols. Reject if the syndrome differs from all $q$-bit syndromes of $\mathcal S$. There are at most $(1+M2^j)^{k+1}$ such vectors, so $q=O(k\log k)$. After composition with the derivative maps, each syndrome bit uses one parity query to $f$. If at most $k$ raw variables are essential, this test always passes. If at least $k+1$ are essential, a fixed set of $k+1$ of them lands in different buckets with probability above $99/100$. All of the corresponding symbols remain nonzero with probability above $99/100$. Whenever $V$ has more than $k$ nonzero symbols, its independent parity sketch collides with $\mathcal S$ with probability at most $1/128$.

For the second test identify every variable in a bucket with one new variable, giving an $M$-variable function $g$. Any parity of the truth table of $g$ is a parity of $f$, by collecting the signed multiplicities modulo two (all substitutions here use positive variables). A size-$s$ circuit for $f$ gives one for $g$, so yes inputs still pass with certainty. If $f$ has at most $k$ essential axes, then with probability above $99/100$ they hash injectively. On that event $g$ and $f$ have the same minimum size, by relabeling and fixing dummy variables. Fingerprint $g$ against all circuits of size $s$ on $M$ variables; with $O(s\log(M+s))=O(s\log s)$ binary parities, a fixed hard $g$ passes with probability at most $1/128$. This proves the upper bound in all cases. The number of addresses in a parity query and the computation between queries are unrestricted. Earlier parity decision tree work treats the related Hamming-weight and $k$-linearity testing problems \cite{bsh}.

\begin{thebibliography}{9}
\bibitem{intmult} D. Harvey and J. van der Hoeven. \href{https://doi.org/10.4007/annals.2021.193.2.4}{Integer multiplication in time $O(n\log n)$}. Annals of Mathematics, 2021. A quasi-linear Boolean circuit bound suffices here.
\bibitem{polyprod} D. Cantor and E. Kaltofen. \href{https://kaltofen.math.ncsu.edu/bibliography/91/CaKa91.pdf}{On fast multiplication of polynomials over arbitrary algebras}. Acta Informatica, 1991.
\bibitem{bsh} \href{https://eccc.weizmann.ac.il/report/2013/142/}{Property testing bounds for linear and quadratic functions via parity decision trees}. ECCC TR13-142, 2013.
\bibitem{hl} W. M. Hoza and Z. Lv. \href{https://doi.org/10.4230/LIPIcs.APPROX/RANDOM.2025.35}{Fooling near-maximal decision trees}. RANDOM 2025.
\bibitem{cher} M. Cheraghchi. \href{https://arxiv.org/abs/1107.4709}{Applications of derandomization theory in coding}. Doctoral dissertation, 2010.
\bibitem{chash} A. V. Chashkin. \href{https://keldysh.ru/papers/2024/mvk/mvk2024_152.pdf}{On computing partial Boolean functions}, Theorem 6.1, 2024.
\bibitem{lup} O. B. Lupanov. A method of circuit synthesis. Izvestiya VUZ, Radiofizika 1 (1958), 120--140.
\bibitem{hk} W. Hoza and A. Klivans. \href{https://eccc.weizmann.ac.il/report/2016/172/}{Preserving randomness for adaptive algorithms}. ECCC TR16-172, revised version.
\bibitem{fg} M. A. Forbes and V. Guruswami. \href{https://arxiv.org/abs/1411.7455}{Dimension expanders via rank condensers}. APPROX/RANDOM, 2015.
\bibitem{rank} S. Ballet and J. Pieltant. \href{https://arxiv.org/abs/1003.1864}{On the tensor rank of multiplication in any extension of $\mathbb F_2$}. Finite Fields and Their Applications, 2011.
\end{thebibliography}
\end{document}
